\chapter{In the field}\label{ch:field}
\begin{goals}
\textbf{In this chapter} you measure planes and lines with the phone, refer them to true north, attach
outcrop photos, and publish the log, signed, for the desk.
\end{goals}

On a phone, wuffle reads the device's orientation and turns it into an attitude. The net fills as you
work, so you see the fabric take shape at the outcrop instead of that evening.

\shot[0.5\linewidth]{field-net}{The compass: the mode switch, the live reading over the net, the measure
button, and the log below.}

\section{The screen}\label{sec:field-screen}

From the top:

\begin{description}
\item[Mode.] \ui{plane} or \ui{line}: what the next measurement is.
\item[The strip.] The position's accuracy (\ui{±4 m}) and which north the readings use
(\ui{true N · decl −15.2°}).
\item[The net.] An equal-area, lower-hemisphere net. The live reading is a dashed ghost: a great circle
for a plane, a point for a line. In the corners: the reading as numbers (dip direction and dip, or trend
and plunge), a bar that fills with the dip, and the measure button \ui{◎}.
\item[Typed entry.] A field for measurements you type (\cref{sec:manual-entry}).
\item[The log.] The count, \ui{photo}, \ui{clear} and \ui{publish →}, then every measurement, newest
first. A plane shows its strike too. The \ui{×} on a row deletes it.
\end{description}

\section{Measuring}\label{sec:measuring}

For a \textbf{plane}, lay the phone flat on the surface, screen up. wuffle reads the plane the phone lies
in: its dip direction and dip.

For a \textbf{line}, choose \ui{line} and lay the phone's long edge along the lineation, top edge pointing
down the plunge. The trend and plunge are those of the phone's long axis, toward its top.

Tap \ui{◎} to keep the reading. The phone buzzes, the measurement is drawn in place of the ghost and it
joins the log.

\begin{tip}
Inside the lead-acid shell, wuffle reads the phone's fused orientation sensor, which also reports how
well the magnetometer is calibrated. When it says \ui{compass low} or \ui{unreliable}, wave the phone in a
figure of eight until it says \ui{calibrated}, then measure.
\end{tip}

\section{True north}\label{sec:true-north}

A phone's compass points to magnetic north. wuffle corrects each reading to \term{true north} with the
\term{declination} at your position, from the World Magnetic Model (WMM2025) that ships inside it. It
needs a position for that, so the correction starts with the first GNSS fix:

\begin{itemize}
\item before a fix, readings are magnetic, and the strip shows no declination;
\item after it, they are true, and the strip shows the declination in use. It is recomputed when you have
moved about a tenth of a degree (some 10~km).
\end{itemize}

Every measurement records which it was. The log's \col{north} column says \code{true} or
\code{magnetic}, and \col{declination\_deg} gives the correction applied. A typed measurement says
\code{as-typed}: it is referred to whatever north you used when you wrote it down.

\begin{caution}
Outside the shell (wuffle in an ordinary phone browser) there is no position, so readings stay magnetic.
The \col{north} column says so on every row; correct them before comparing with true-north data.
\end{caution}

\section{Typing measurements}\label{sec:manual-entry}

The field under the net takes typed measurements: \code{045/30} is a plane, \code{310→70} a line, and
\code{310>70} is the same line without hunting for the arrow. You can also type \code{310-70 line}; and
with the \ui{line} mode on, \code{310/70} is taken as a line.
Press \key{Enter} or \ui{add}. This is how the compass works on a computer, and how you correct a reading
in the field: delete it and type the right one.

\section{Photos}\label{sec:photos}

\ui{photo} opens the camera. \ui{capture} takes a picture and attaches it to the \emph{last}
measurement: the row gets a camera mark and the log's \col{photo} column names the file
(\file{wuffle-002-2026-11-20T140540.jpg}: the measurement's number and time). Measure first, then
photograph.

\section{Publishing the log}\label{sec:publish}

\ui{publish →} writes the log out:

\begin{itemize}
\item \file{wuffle-log.csv} to \file{Downloads}: one row per measurement, with its kind, attitude,
strike, time, position and accuracy, north, declination and photo;
\item \file{wuffle-log.csv.sig} beside it, the phone's \term{signature} over the exact bytes of the CSV
(inside the shell, when the phone has a hardware key);
\item each photo to \file{Pictures}.
\end{itemize}

The signature is made by a key that lives in the phone's secure hardware and never leaves it. It lets
anyone check later that the CSV is byte for byte what this phone wrote, and the desk does exactly that
when it opens the log (\cref{sec:field-logs}).

On a computer, \ui{publish →} downloads the CSV and the photos; there is no hardware key there, so there
is no signature.

\begin{tryit}
\begin{enumerate}
\item On a computer, open wuffle with \code{?field} at the end of the address.
\item Type \code{120/45}, \key{Enter}; \code{118/47}, \key{Enter}; \code{200>20}, \key{Enter}. Two
planes and a line are drawn.
\item Press \ui{publish →}. \file{wuffle-log.csv} downloads.
\item Open it on the desk (\menu{File>Open data or a project…}). The planes and the line become two
datasets, \code{field … · planes} and \code{field … · lines}, named for the log's day, and the CSV
itself is kept as it came (\cref{sec:field-logs}).
\end{enumerate}
\end{tryit}

\section{Taking a log back}

Inside the shell, a CSV shared to wuffle (or opened with it from a file manager) is read back into the
log, as long as it has the log's \col{type}, \col{dip\_dir\_or\_trend} and \col{dip\_or\_plunge}
columns. That is how you carry on a log another day, or on another phone.
